\section{Conclusion}
\label{sec:conclusion}

The aim of this project, set in Chapter \ref{chap:new-ideas}, was to design, implement, and evaluate the performance of a "prototype access control system that uses \ac{gr} to passively identify subjects".

Referring to the results of Chapter \ref{chap:implementation} --- Implementation and Chapter \ref{chap:results-and-evaluation} --- Results and Discussion, the project can be considered a success, both functional, and personal in respect to the researcher.

\subsection{Achievements}

This project successfully achieved the core aim outlined in the \ac{ppd} and refined in Chapters \ref{chap:context} and \ref{chap:new-ideas}, with all five objectives met as well as all mandatory and desired functional and non-functional requirements.
The system's architecture was implemented as planned, with a client-server topology communicating via an \ac{api}.
This supports scalability, with the potential for a multi-client-server setup to be implemented in the future --- beneficial in the context of a \ac{pacs}.

Significant technical achievements of this project include the implementation of a \ac{cnn} + \ac{pca} + cosine similarity classificaiton pipeline for identifying users based on a single image representation of their statistical and spatio-temporal gait features --- a \ac{gei}.
Fully functional identity registration and classification pipelines were also implemented,which enable live and asynchronous demonstration of the systems capabilities, aiding in the justification of implementing such a system in a real-world setting. The manual classification feature and \ac{grs} 'register' and 'classify' endpoints  significantly supporting the evaluation of the system's performance and accuracy.

All of this was achieved while using medium-powered consumer-grade hardware (defined in Section \ref{sec:technologies-used-and-techniques-applied}) for development and evaluation. 
This proves the system's efficiency and feasibility for cost-effective implementation in a real-world setting, especially when compared to the assumed cost of the only existing example of a \ac{gr}-based \ac{pacs} system --- \cite{stepscanGaitAuthenticationSecurity2020}.

\subsection{Criticism}

Despite the overwhelming success of the project, several issues were encountered, throughout both development and evaluation. 
The system also faces several limitations that have not been mitigated, these issues and limitations include:

\paragraph{Overfitting} The initial choice of classification approach was a $k$-\ac{nn} model. 
Unfortunately, this displayed high sensitivity when tested on small datasets, resulting in overfitting where subjects were commonly misclassified, and confidently so.
This necessitated futher development time to transition to a \ac{cnn} + \ac{pca} + cosine similarity based approach, which has proven much more effective, as evidenced by the results in Chapter \ref{chap:results-and-evaluation}.
There still remains room for improvement, with a more targeted approach --- the pre-trained \codeword{EfficientNetB0} model is a generic image classification model, and could be replaced with a model trained on a gait-specific dataset.

\paragraph{Data Sensitivity and Environment Inconsistency} Although considered a success, the effectiveness of the silhouette extraction process is dependent on the contrast between the subject and background, meaning in varied environments, and with different subjects and clothing, accuracy is reduced.

\paragraph{Test Diversity} While performance testing was conducted using a secondary dataset of over 100 subjects, primary testing only consisted of six (6) subjects, albeit with a balanced sex ratio: 3 females and 3 males.
The small sample size may have resulted in potential issues being missed, meaning the results may not represent the performance of the system in a real-world setting with far more subjects.

\paragraph{Ethical Concerns} Most negative-leaning feedback received from the initial public opinion survey (see Section \ref{app:public-opinion-survey}) and post-user testing questionnaire (see Section \ref{subsec:user-testing}) involved concerns about the potential misuse --- especially relating to surveillance --- of \ac{gr}-based identification technology.
If adopted in a real-world setting, legislation and policies must be updated to control the use of behavioural biometric factors, like how existing legislation controls the use of \ac{pii}.


\subsection{Future Work}

As it stands, several improvements could be made to this system, these have been briefly covered throughout the report, but include the following:

\paragraph{Multiview Capture} As discussed in Chapter \ref{chap:context}, state-of-the-art proposals in \ac{gr} utilise multiple cameras to capture different angles of subjects. 
Increasing the number of datapoints available would likely have a direct impact on classification precision and robustness.
OpenPose, the model-based pose estimation framework discussed in Section \ref{subsubsec:model-based-approaches} supports multi-view input, infact it is required if wanting to produce a 3-dimensional pose estimate.
Given suitable hardware is available, an 3D model-based approach may be a more effective demonstration of the potential accuracy of \ac{gr}.
Alternatively, multi-view \ac{gei} could be implemented, which would likely still have an positive impact on accuracy.

\paragraph{Hardware Optimization} Although development and evaluation was performed on a consumer-grade laptop, future research could focus on optimising the system for deployment on low(er)-powered edge devices (e.g., the Nvidia Jetson Nano). If successful, this may increase interest in adopting \ac{gr} in a real-world setting due to cost-effectiveness.

\paragraph{Accessibility Enhancements} A requirement that was categorised using the MoSCoW prioritisation technique as a 'Won't Have' was to implement back-up access methods as a failsafe for when \ac{gr} is not possible due to environmental conditions, or subject impairment (e.g., wheelchair users).
Implementing alternative access methods and producing a multi-modal access control architecture would be more comprehensive.

\subsection{Final Thoughts}

This project has validated the sanity in a proposition for a \ac{gr}-based \ac{pacs}, enabling user-independent authentication by leveraging behavioural biometrics to assess subjects from a distance.

Although only a \ac{poc}, the system goes far beyong what would be considered an \ac{mvp} and is fully functional, as proven by demonstration of it's capabilities to volunteers and during functional and performance testing.

Hopefully future research and development will be inspired by these findings, with the public sentiment and acceptance of \ac{gr} increasing.